\section{Surrender Before Evil and Suffering}

\subsection{Evil is not one thing}

The word ``evil'' can name different privations. Moral evil is the free creature's failure of due order to reason, divine law, and the last end. Physical or suffered evil is the loss of a created good: health, safety, reputation, property, relationship, bodily integrity, or life. A just penalty is again distinct, for its privation can belong to an order of justice. These senses must not be collapsed into the claim that whoever suffers is being punished for a personal fault.

Christ rejects that inference. Neither the man born blind nor his parents sinned in the manner his questioners presumed (John 9:1--3); the Galileans killed by Pilate were not worse sinners because they suffered, nor were those crushed by the tower (Luke 13:1--5). All creation groans under corruption while awaiting redemption (Rom 8:18--25). Some suffering follows one's sin, some another's sin, some the finite interaction of natural causes, and much exceeds our power to assign a particular reason.

God is never the cause of the moral defect. He is Lord of every created good lost, every natural power exercised, every limit on the offender, every judgment, every possibility of rescue, and the final history. ``Permission'' does not mean that God stands temporarily absent. It means that his wise positive will gives no evil its evil, even while his sustaining and ordering causality encompasses the world in which the defect occurs.

\subsection{One event, four moral relations}

When Joseph is sold, at least four relations must be distinguished.

\begin{enumerate}[leftmargin=1.6em]
\item \textbf{The brothers' relation:} they choose betrayal from envy and remain guilty.
\item \textbf{Joseph's relation:} he suffers an injustice he did not choose; later he acts prudently, tests, forgives, and provides.
\item \textbf{The created causal history:} merchants, famine, political office, dreams, storage, travel, and recognition really contribute.
\item \textbf{God's providential relation:} without sharing the evil intention, God orders the whole toward preservation and reconciliation: ``you meant evil against me, but God meant it for good'' (Gen 50:20).
\end{enumerate}

The grammar reaches its summit in the Passion. The Father gives the Son; the Son gives himself; Judas hands him over; authorities condemn; soldiers execute; disciples flee; the women remain; Joseph of Arimathea buries. The external history belongs to one saving mystery, but the moral acts are not one act. Augustine notices the decisive point: Father, Son, and Judas can be said to ``hand over'' Christ, yet charity makes the Father's and Son's self-gift saving while avarice makes Judas's betrayal wicked (\emph{Tractates on First John} VII.7). The Resurrection vindicates Providence without retroactively blessing treachery.

\subsection{Acceptance does not ratify the wrongdoer}

Acceptance is not a feeling, instant psychological integration, or consent to ongoing danger. It is the gradual, graced refusal to make the impossible undoing of the past a condition of remaining with God, while truth and charity still govern every present response. It does not say that the event's moral defect was good, that the human cause does not matter, or that future harm should be allowed.

\begin{threestable}{Surrender says}{Surrender does not say}{Charity may require}
This has happened within a world still governed by God. & God committed, approved, or commanded the sin. & Name the wrong precisely; preserve evidence; report it.\\
No wound can force me outside the possibility of grace and final good. & The wound is small, deserved, or emotionally easy. & Lament; seek sacramental, medical, psychological, legal, or pastoral help.\\
I relinquish vengeance and final judgment to God. & Justice, boundaries, restitution, and consequences are unchristian. & Protect another; leave danger; seek a restraining or judicial remedy.\\
I ask for the offender's conversion and refuse hatred. & Reconciliation requires immediate access, trust, or restored office. & Forgive as a will toward the person's true good while prudence governs contact.\\
I will obey God in the next possible act. & I must discover the secret reason God allowed this. & Act on known duties and tolerate not knowing the hidden design.\\
\end{threestable}

This distinction is especially necessary when power is unequal. Advice about surrender given to a victim but not repentance and restitution to an offender becomes a weapon. Catholic surrender never transfers the sinner's culpability to the sufferer. It never makes another's access to my body, children, home, money, office, or confidence a proof of forgiveness.

\subsection{The limits of knowing why}

Saint-Jure and La Colombière repeatedly invite the reader to trust that trials lie within a wise and useful order. This can be medicine against the fantasy of a pointless universe. It becomes presumption if one claims to know the particular temporal benefit God selected in every tragedy. Job's friends speak too confidently about reasons; God answers from a wisdom larger than their formulas. The \emph{Catechism} says the ways of Providence are often unknown until the end (CCC 314).

The safe claim is both stronger and humbler: no created evil becomes good, no evil escapes God's permission and limit, and no evil can prevent omnipotent charity from attaining the ultimate good God wills. Romans 8:28 is addressed to those who love God and are called according to his purpose; it is a promise of divine operation toward final good, not a demand that the sufferer call every isolated event beneficial on sight.

Nor does 1 Corinthians 10:13 promise that unaided psychological powers will never be overwhelmed by any trauma. It speaks of temptation and of God's faithful provision of an escape so that one may endure. The verse therefore may not be used to tell traumatized people that their suffering lay within unaided psychological capacity. In danger, another person, medicine, rest, a locked door, a court order, the police, a therapist, a confessor, or flight may be ordinary secondary causes Providence gives for safety; that pastoral claim is not presented as the verse's direct exegesis.

\subsection{Against fatalism and Quietism}

Fatalism makes outcomes follow from blind or impersonal necessity. Providence is the personal order of wisdom and love that includes contingent and free causes. Quietism, in its condemned forms, treated spiritual passivity as perfection: petition, active resistance to temptation, concern for virtue, and deliberate acts could be dismissed as obstacles. Innocent XI's \emph{Caelestis Pastor} (1687; DS 2201--2269) rejected propositions that made the soul inanimate, unconcerned with salvation and virtue, unwilling to petition, or passive before temptation.

The Catholic opposite is not anxious activism. It is cooperation. Prudence discerns the good and the means (CCC 1806). Prayer is a gift and a determined response that entails effort (CCC 2725). Filial trust perseveres in asking (CCC 2734--2741). Christian life joins prayer to works and works to prayer (CCC 2742--2745).

The devotional anthology, read as a whole, belongs on this active side. It places commandments and duties first, permits reasonable remedies and counsel, devotes a major section to persistent prayer, and describes surrender as doing what lies within one's power before yielding the result. Its occasional command to deny ``discussion'' can safely be received here as relinquishing sterile counterfactual rumination after prudent judgment, not surrendering reason, conscience, or safeguards.

\begin{studybox}{The test of authentic abandonment}
If ``surrender'' makes truth less truthful, duty less binding, prayer less urgent, the vulnerable less protected, repentance less concrete, or charity less active, it is counterfeit. Providence gives secondary causes their dignity; trust uses them without worshiping their success.
\end{studybox}
